\documentclass[11pt,a4paper]{article}

\usepackage[T1]{fontenc}
\usepackage[utf8]{inputenc}
\usepackage{lmodern}
\usepackage[a4paper,margin=26mm]{geometry}
\usepackage{amsmath,amssymb,amsthm,mathtools}
\usepackage{booktabs,array}
\usepackage{enumitem}
\usepackage{microtype}
\usepackage{xcolor}
\usepackage{url}
\usepackage[colorlinks=true,linkcolor=blue!55!black,citecolor=blue!55!black,
            urlcolor=blue!55!black]{hyperref}

\hypersetup{
  colorlinks=true,
  linkcolor=blue!55!black,
  citecolor=blue!55!black,
  urlcolor=blue!55!black,
  pdftitle={Message Retargeting in the NGCC CompactSQIsign2D2 Submission},
  pdfauthor={Martin Feussner}
}

\newtheorem{theorem}{Theorem}
\newtheorem{corollary}[theorem]{Corollary}
\newtheorem{remark}[theorem]{Remark}
\newcommand{\CSQI}{\mathsf{CompactSQIsign2D2}}
\newcommand{\Hash}{\mathsf{HASH}}
\newcommand{\Rec}{\mathsf{Recover}}
\newcommand{\Verify}{\mathsf{Verify}}
\newcommand{\Sign}{\mathsf{Sign}}
\newcommand{\KeyGen}{\mathsf{KeyGen}}
\newcommand{\Codom}{\operatorname{Codom}}

\title{\bfseries Message Retargeting in the NGCC
CompactSQIsign2D2 Submission\\[0.25em]
\large A One-Query Fresh-Message Forgery Across All Eight Parameter Sets}
\author{Martin Feussner\\
\small Selmer Center, University of Bergen\\
\small \texttt{martin.feussner@uib.no}}
\date{28 September 2026}

\begin{document}
\maketitle

\begin{center}
\fbox{\parbox{0.92\linewidth}{\small\textbf{Provenance and review status.}
The attack was independently found by an AI system, \mbox{OpenAI Codex} using
Daybreak Blue at maximum reasoning effort, without a human first proposing
this specific signature attack.  The AI also produced the derivation,
implementation, experiments, and initial manuscript.  At the date above, no
human cryptanalyst has independently verified the argument, implementation,
measurements, or novelty assessment.  This preliminary disclosure is provided
for human review and reproduction.}}
\end{center}
\vspace{0.7em}

\begin{abstract}
The compact mode in the NGCC SQIsign2D2 submission serializes both a description of the dual
challenge isogeny and a Fiat--Shamir challenge scalar.  Its verifier reconstructs a challenge
curve from the response, follows the transmitted dual path back to a commitment curve, hashes
that commitment with the public key and message, and compares the result with the transmitted
scalar.  The scalar never constrains either isogeny path.  In particular, verification omits the
condition that the forward isogeny selected by the hash-derived scalar has the reconstructed
challenge curve as its codomain.

This omission gives a deterministic one-query EUF-CMA forgery against the submitted compact
mode.  From one valid signature, an attacker publicly reconstructs the same commitment curve,
hashes it with any fresh target message, and replaces only the serialized challenge field.  Every
geometric check receives its original input, while the final hash comparison now holds for the
target message.  No collision, second preimage, secret isogeny, or equivalent key is required.

An independent implementation of this transform against the unmodified reference sources
produced 24 accepted fresh-message forgeries: three for each of the eight submitted compact
parameter sets.  Negative controls rejected the untouched signature on each target and the
retargeted signature on its source message; byte comparisons found no change outside the
challenge field.  Median forgery times ranged from 20 milliseconds to 1.98 seconds on an Intel
Core i5-6500 laptop.  The flaw is present in the NGCC verification algorithm and implementation.
It is not a claim against the original 2025 CompactSQIsign2D2 construction, whose verification
algorithm includes a separate equality binding the hash-derived forward kernel to the transmitted
dual path.
\end{abstract}

\section{Introduction}

SQIsign2D2 is an isogeny-based Fiat--Shamir signature submitted to the
Next-generation Commercial Cryptographic Algorithms Program in both uncompressed and compact
modes~\cite{sqisign2d2-spec}.  At a high level, the signer
commits to a curve $E_{\rm com}$, derives a scalar
\[
  c=\Hash(\mathit{pk}\parallel j(E_{\rm com})\parallel m),
\]
and uses $c$ to select a $3^b$-isogeny from $E_{\rm com}$ to a challenge curve
$E_{\rm chl}$.  A response isogeny connects the public-key curve to the same
$E_{\rm chl}$.  Sound Fiat--Shamir verification must bind these three objects: the response,
the commitment, and the hash-derived challenge~\cite{fiat-shamir}.

The NGCC compact mode reconstructs rather than transmits the commitment curve.  Its signature
contains data from which the verifier obtains $E_{\rm chl}$ and a scalar description of the
dual challenge isogeny
\[
  \widehat\phi:E_{\rm chl}\longrightarrow E_{\rm com}.
\]
It also serializes $c$.  The verifier follows $\widehat\phi$ to obtain $E_{\rm com}$ and checks
that the serialized $c$ equals the hash of the message and this curve.  It does not verify that
$c$ selects the forward isogeny dual to $\widehat\phi$.

This distinction is decisive.  In the submitted verification relation, the response and dual-path data are independent of the
serialized field at verification time.  Once a valid tuple has passed all geometric conditions,
the attacker may recompute only the detached hash field for a different message.  The same
response then authenticates every target message.

The attack has the following properties.
\begin{enumerate}[leftmargin=1.7em]
  \item It is a fresh-message forgery in the standard EUF-CMA game, using one signing query.
  \item It succeeds deterministically for any fresh target message and every submitted
        compact parameter set.
  \item It requires only computations already performed by public verification.
  \item It changes one contiguous field of a valid serialized signature and leaves the full
        isogeny response unchanged.
  \item It does not rely on malformed curves, unchecked inputs, assertion behavior, or any
        previously reported implementation issue.
\end{enumerate}

Section~\ref{sec:construction} isolates the relevant part of compact signing and verification.
Section~\ref{sec:attack} gives the transform and proves the forgery.  Section~\ref{sec:experiment}
reports validation against all eight reference implementations.  Section~\ref{sec:scope}
distinguishes the NGCC submission from the original published compact construction and from
other public findings, and
Section~\ref{sec:repair} gives a repair condition.

\section{The compact construction}
\label{sec:construction}

We use only the interfaces needed for the attack.  The internal quaternion and
two-dimensional-isogeny computations are irrelevant.

\subsection{Commitment and forward challenge}

The signer samples a commitment curve $E_{\rm com}$ and obtains a canonical basis
$(R_{\rm com},S_{\rm com})$ of its $3^b$-torsion.  It computes
\begin{equation}
  c=\Hash(\mathit{pk}\parallel j(E_{\rm com})\parallel m)
     \in \mathbb Z/3^b\mathbb Z                                      \label{eq:hash}
\end{equation}
and the challenge isogeny
\begin{equation}
  \phi_c:E_{\rm com}\longrightarrow E_{\rm chl},\qquad
  \ker(\phi_c)=\langle R_{\rm com}+[c]S_{\rm com}\rangle .          \label{eq:forward}
\end{equation}
This is Algorithm~3.5, lines 14--15, of the submission specification.  The response is then
constructed so that its relevant codomain is the same curve $E_{\rm chl}$.

For compact encoding, the signer computes the dual path
\begin{equation}
  \widehat\phi_c:E_{\rm chl}\longrightarrow E_{\rm com}.             \label{eq:dual}
\end{equation}
Relative to a canonical $3^b$-torsion basis $(R_{\rm chl},S_{\rm chl})$, its kernel is encoded
by a projective scalar $t$ and a selector $g$:
\[
  \ker(\widehat\phi_c)=
  \begin{cases}
    \langle R_{\rm chl}+[t]S_{\rm chl}\rangle,&g=0,\\
    \langle [t]R_{\rm chl}+S_{\rm chl}\rangle,&g=1.
  \end{cases}
\]
The remaining compact data, denoted below by $\rho$, describes the response and supplies the
hints and matrices from which the verifier recovers $E_{\rm chl}$.

The submitted byte encoding is therefore conveniently written
\begin{equation}
  \sigma=(\rho,t,g,c).                                                \label{eq:sig}
\end{equation}
Here $\rho$ includes the transmitted curve, response matrix, curve-selection bit, and basis
hints.  The exact names of these fields do not affect the attack.

\subsection{Compact verification}

Algorithm~3.6 of the NGCC specification and the reference implementation perform the following operations.
\begin{enumerate}[leftmargin=1.7em]
  \item From $\mathit{pk}$ and $\rho$, compute a dimension-two isogeny and select its indicated
        elliptic factor as $E_{\rm chl}$.
  \item From $E_{\rm chl}$ and $(t,g)$, evaluate the $3^b$-isogeny in
        Equation~\eqref{eq:dual} and call its codomain $E_{\rm com}$.
  \item Compute $c^*=\Hash(\mathit{pk}\parallel j(E_{\rm com})\parallel m)$ and test
        $c=c^*$.
  \item Perform any remaining validity and common-part checks on the response and the encoded
        dual path.
\end{enumerate}

Let
\begin{equation}
  (E_{\rm chl},E_{\rm com})=\Rec(\mathit{pk},\rho,t,g)                \label{eq:recover}
\end{equation}
denote the first two steps.  Crucially, $\Rec$ does not take $c$ or $m$ as input.  All remaining
geometric checks are also independent of $c$ and $m$.  Consequently there is a predicate
$P$ such that the implemented acceptance condition has the form
\begin{equation}
  \Verify(\mathit{pk},m,(\rho,t,g,c))=1
  \quad\Longleftrightarrow\quad
  P(\mathit{pk},\rho,t,g)=1
  \;\wedge\;
  c=\Hash(\mathit{pk}\parallel j(E_{\rm com})\parallel m).           \label{eq:accept}
\end{equation}

Equation~\eqref{eq:accept} checks that $c$ is a digest of the reconstructed commitment, but
it does not use $c$ to check the challenge isogeny.  The omitted condition is
\begin{equation}
  E_{\rm chl}\ \cong\
  \Codom\!\left(E_{\rm com}/
  \langle R_{\rm com}+[c]S_{\rm com}\rangle\right),                  \label{eq:missing}
\end{equation}
together with compatibility between that forward isogeny and the encoded dual path.  The
honest signer creates this relation, but the verifier never tests it.

\subsection{A specification tuple inconsistency}
\label{sec:tuple}

The NGCC specification contains a tuple inconsistency about $c$.  Table~1.5 states that a
compact signature contains the challenge.  Algorithm~3.6, line 20, compares a variable
\texttt{chl} with the hash, and Table~6.1 reports sizes that include a full challenge field.
The reference encoding likewise serializes that field.  On the other hand, Algorithm~3.5,
line 47, and the input tuple printed for Algorithm~3.6 omit it.

If those tuple omissions are read literally, verification refers to an undefined value and the
scheme is incomplete.  The coherent interpretation supported by the notation table, size table,
KAT modes, and executable submission is Equation~\eqref{eq:sig}: $c$ is transmitted and
compared with the hash.  The attack applies to that submitted interpretation.  Adding $c$ to the
printed tuples resolves the syntax but does not add the missing relation
\eqref{eq:missing}.

\section{Universal message retargeting}
\label{sec:attack}

Let $(\mathit{pk},\mathit{sk})\leftarrow\KeyGen$ and request a signature on any message
$m_0$:
\[
  \sigma_0=(\rho,t,g,c_0)\leftarrow\Sign(\mathit{sk},m_0).
\]
For any fresh target message $m_1\ne m_0$, perform the following public computation.
\begin{enumerate}[leftmargin=1.7em]
  \item Compute $(E_{\rm chl},E_{\rm com})=\Rec(\mathit{pk},\rho,t,g)$ exactly as the verifier.
  \item Compute
        \[
          c_1=\Hash(\mathit{pk}\parallel j(E_{\rm com})\parallel m_1).
        \]
  \item Output
        \[
          \sigma_1=(\rho,t,g,c_1).
        \]
\end{enumerate}
Only the challenge field is replaced.  The attacker does not modify either isogeny path.

\begin{theorem}[One-query universal forgery for the submitted compact mode]
Assume honest compact signatures are accepted.  The transform above wins the EUF-CMA game
against the NGCC compact mode after one signing query.  It succeeds for every fresh target
message $m_1$ and requires one public commitment reconstruction and one evaluation of $\Hash$.
\end{theorem}

\begin{proof}
Because $\sigma_0$ is accepted, Equation~\eqref{eq:accept} gives
$P(\mathit{pk},\rho,t,g)=1$.  The forgery keeps every input to $P$ unchanged, so the same
predicate remains true.  Recovery from $\sigma_1$ also receives exactly the original
$(\rho,t,g)$ and therefore returns the same $E_{\rm chl}$ and $E_{\rm com}$.  By construction,
the new field satisfies
\[
  c_1=\Hash(\mathit{pk}\parallel j(E_{\rm com})\parallel m_1).
\]
Both conjuncts in Equation~\eqref{eq:accept} hold, hence
$\Verify(\mathit{pk},m_1,\sigma_1)=1$.  The message $m_1$ was never submitted to the signing
oracle, so $(m_1,\sigma_1)$ is a valid EUF-CMA forgery.
\end{proof}

\begin{corollary}
One submitted compact signature can be retargeted independently to arbitrarily many fresh messages.
Neither the size of the challenge space nor the number of iterated SM3 calls changes the
success probability.
\end{corollary}

The result does not ask for a collision.  The attacker evaluates the hash honestly for each
new message and writes its output into a field that verification treats as independent data.
Likewise, the attacker does not solve an isogeny problem: $\Rec$ is the public algorithm needed
for ordinary verification.  The hash input order used by the reference code differs cosmetically
from the displayed specification order, but the attacker calls the same submitted hash routine,
so either ordering admits the transform.

\section{Reference-code validation}
\label{sec:experiment}

\subsection{Source-level data flow}

The compact branch of the Level1-eff reference verifier illustrates the common code path.
Lines 1136--1186 of \texttt{sqisign/sign.c} recover and standardize $E_{\rm chl}$ using the
public key and non-challenge signature fields.  Lines 1188--1201 derive the dual kernel solely
from \texttt{sig->ind}, \texttt{sig->scalar}, and the basis hints, then evaluate it to obtain
$E_{\rm com}$.  Only after this reconstruction do lines 1203--1205 read
\texttt{sig->chall}, comparing it with the new hash.  No earlier or later operation uses that
field.

The serialized layout makes the transformation particularly simple.  In compact mode,
\texttt{pack.c} writes the encoded curve first and the challenge immediately afterward
(lines 232--235); decoding uses the same order (lines 301--305).  The attacker replaces the
contiguous byte interval
\begin{equation}
 [\,\texttt{FP2\_ENCODED\_BYTES},\;
    \texttt{FP2\_ENCODED\_BYTES}+\texttt{NWORDS\_ORDER\_3\_BYTES}\,)
                                                                    \label{eq:bytes}
\end{equation}
and copies every other byte verbatim.

The seven non-Level1-eff instances have byte-identical \texttt{sign.c} files.  The Level1-eff
file differs only in unrelated source details and has the same data flow.  The attack harness
links the original key generation, signing, parsing, hashing, and verification sources without
editing them.  Its attacker-side routine copies the verifier's public reconstruction through the
point where $E_{\rm com}$ is available, invokes the submission's
\texttt{hash\_to\_challenge} routine, and
reserializes the signature.

\subsection{Method and controls}

For each of the eight submitted parameter sets we compiled the portable reference sources with
\texttt{-O2 -DNDEBUG -DCOMPRESSED=1}, GCC 13.3.0, and GMP 6.3.0.  The independent
reproduction ran on 64-bit Linux on a four-core Intel Core i5-6500 processor.  It generated one
key pair and one signature on $m_0$, then retargeted that signature to three distinct fresh
messages.  The same public transform and controls were fixed for every parameter set.

Each trial checked all of the following:
\begin{enumerate}[leftmargin=1.7em]
  \item the honest signature is accepted on $m_0$;
  \item the untouched signature is rejected on a fresh message;
  \item the retargeted signature is accepted on its fresh message;
  \item the retargeted signature is rejected on $m_0$; and
  \item all bytes outside the interval in Equation~\eqref{eq:bytes} are unchanged.
\end{enumerate}
The reported forgery time includes signature parsing, both public isogeny reconstructions, all
configured hash grinding, replacement, and serialization.  It excludes key generation and the
single oracle signature.  Times are the median and range of the three fresh targets.

\begin{table}[t]
\centering
\small
\caption{End-to-end compact-signature retargeting.  Every fresh forgery was accepted.}
\label{tab:results}
\begin{tabular}{@{}lrrrrrr@{}}
\toprule
Parameter set & $b$ & Signature & Offset & Field & Median & Range \\
              &     & bytes     & bytes  & bytes & ms     & ms \\
\midrule
Level1-eff &  78 & 168 &  64 & 16 &   63 &   56--67 \\
Level1-sec &  84 & 178 &  68 & 17 &   71 &   70--88 \\
Level2-eff &  96 & 208 &  80 & 20 &   20 &   19--20 \\
Level2-sec & 101 & 218 &  84 & 21 &  161 & 128--163 \\
Level3-eff & 156 & 326 & 128 & 31 &  340 & 334--346 \\
Level3-sec & 163 & 338 & 132 & 33 &  456 & 344--478 \\
Level5-eff & 319 & 648 & 256 & 64 & 1978 & 957--2220 \\
Level5-sec & 334 & 678 & 268 & 67 & 1689 & 668--1869 \\
\bottomrule
\end{tabular}
\end{table}

All 24 fresh-message forgeries passed every check.  Every measured transform completed within
2.23 seconds.  The controls establish that acceptance was caused by the new challenge field,
while the byte comparison establishes that no response component was altered.  The measured
key-generation and one-oracle-signature times are excluded from the attack timings; the signing
query is supplied by the EUF-CMA oracle.

\begin{samepage}
For artifact identification, the tested specification PDF has SHA-256 digest
\begin{center}
\small\ttfamily
598139DA24C326890600CED289C41E62\\
1E7DF5F1612EF3C1C699E620E33FB1EC
\end{center}
\end{samepage}

\section{Scope of the result}
\label{sec:scope}

\subsection{The defect is in the submitted compact verification relation}

The proof uses the acceptance structure printed in Algorithm~3.6 of the NGCC submission:
recover
$E_{\rm chl}$ from the response, recover $E_{\rm com}$ from $(t,g)$, and compare a challenge
value with the hash.  Equation~\eqref{eq:missing} is absent from that algorithm.  Thus the
attack is not created by a parser discrepancy or a reference-only shortcut.  It is a flaw in the
verification relation specified and implemented for the submitted compact mode.

The submitted reference verifier performs an additional common-part check after the hash
comparison.  This does not help because the transform leaves the entire response and dual path
unchanged.  Any corrected or strengthened predicate depending only on those original fields has
the same truth value before and after retargeting.  The attack would fail only if verification
made the recomputed challenge constrain the forward challenge isogeny.

\subsection{Distinction from the original published compact construction}

The earlier SQIsign2D2 paper describes a different compact binding mechanism
in Algorithms~8 and 9~\cite{sqisign2d2-eprint}.  Its signature contains a dual-kernel point
$P_2$ and a scalar, denoted $t$ there.  Verification deterministically chooses a complementary
point $Q_2$, evaluates the dual challenge isogeny $\widehat\phi$ with kernel
$\langle P_2\rangle$, sets $L=\widehat\phi(Q_2)$, derives the forward-kernel generator $K_1$
from the message and reconstructed commitment, and requires
\begin{equation}
  K_1=[t]L.                                                        \label{eq:original-binding}
\end{equation}
That equality directly binds the message-derived challenge to the transmitted dual path.

The NGCC submission uses the name $t$ for a projective description of the dual kernel and adds
a transmitted scalar $c$.  It omits Equation~\eqref{eq:original-binding}; its final check is the
detached comparison $c=\Hash(\mathit{pk}\parallel j(E_{\rm com})\parallel m)$.  The present
attack targets that submitted change.  It does not show that Algorithms~8 and 9 of the 2025
paper are vulnerable to the same retargeting transform.

\subsection{Public findings and the uncompressed mode}

As of 28 September 2026, the public NGCC tracker lists two earlier SQIsign2D2 findings:
\texttt{sign-25-1}, a stale-stack verifier failure in the Level2-eff uncompressed reference
implementation, and \texttt{sign-25-2}, a parameter-sizing lead based on improved isogeny
algorithms~\cite{ngcc-tracker}.  Neither finding is the deterministic compact-signature
retargeting described here.

The uncompressed verifier receives $E_{\rm com}$ directly, recomputes the challenge from the
message, and uses that result to form the kernel
$\langle R_{\rm com}+[c]S_{\rm com}\rangle$ before checking the response.  Changing the
message therefore changes an input to the geometric verification.  The detached-field transform
in this note does not apply to that mode.

The present proof assumes an honest, valid compact signature and changes only its challenge
bytes.  It neither uses malformed input nor depends on process state.  It remains valid if all
input validation and common-part checks are implemented as intended, provided the detached
challenge comparison is left unchanged.

\subsection{Security impact}

The EUF-CMA definition in the SQIsign2D2 specification follows the standard adaptive
chosen-message game~\cite{gmr}.  The adversary here asks for one signature on $m_0$ and returns
an accepted pair for a distinct, unqueried $m_1$.  This directly violates the definition.  Since
$m_1$ may be selected after receiving the oracle signature, the failure is stronger than a
same-message encoding malleability result.

The claimed classical security levels of 128, 160, 256, and 512 bits do not enter the attack.
Its cost is comparable to one compact verification and took at most 2.23 seconds in the
submitted portable reference parameter sets tested.

\section{Repair}
\label{sec:repair}

A repair must make the hash output select the same challenge isogeny that the response uses.
Merely recomputing or reserializing $c$ is insufficient.

One option is to restore the binding relation of the original compact construction: transmit the
data needed to define $P_2$ and the binding scalar, reconstruct $L=\widehat\phi(Q_2)$, and check
Equation~\eqref{eq:original-binding} against the hash-derived forward-kernel generator.  The
notation must distinguish that binding scalar from any projective scalar used only to encode the
dual kernel.

Another direct repair for the submitted encoding is:
\begin{enumerate}[leftmargin=1.7em]
  \item reconstruct $E_{\rm chl}$ and $E_{\rm com}$ from the compact response and dual-path
        data;
  \item compute $c=\Hash(\mathit{pk}\parallel j(E_{\rm com})\parallel m)$ internally, without
        trusting a transmitted copy;
  \item derive a canonical basis $(R_{\rm com},S_{\rm com})$ of $E_{\rm com}[3^b]$;
  \item evaluate the forward isogeny with kernel
        $\langle R_{\rm com}+[c]S_{\rm com}\rangle$; and
  \item verify that its standardized codomain and dual kernel agree with the reconstructed
        $E_{\rm chl}$ and the encoded $(t,g)$ path.
\end{enumerate}

At minimum, the codomain equality in Equation~\eqref{eq:missing} must hold.  Checking dual
kernel compatibility as well avoids relying on uniqueness of a path from codomain data alone.
The specification should also state one consistent signature tuple and include this invariant in
the security proof and test suite.  A regression test can take an honest compact signature,
rehash its reconstructed commitment under a different message, replace the challenge field, and
require rejection.

Removing the transmitted $c$ field after deriving it internally would reduce ambiguity, but it
does not by itself provide security: the internally derived value still has to be used in the
forward-path check.

\section{Conclusion}

The NGCC CompactSQIsign2D2 verifier checks a response path and a Fiat--Shamir digest as two
separate facts.  It does not verify the relation that makes the digest the challenge for that path.
This decoupling turns any one valid submitted compact signature into signatures for arbitrary
fresh messages by replacing only the challenge field.  The transform is deterministic, public,
independent of isogeny hardness, and validated against every submitted compact parameter set
in at most 2.23 seconds on the test laptop.  Restoring an explicit forward-to-dual binding, as in
the original published compact construction, prevents this retargeting attack.

\begin{thebibliography}{9}

\bibitem{sqisign2d2-spec}
Zheng Xu, Keteng Huang, Yi Ouyang, Jiankuo Dong, Liang Bai, Yijing Ning,
Wei Wei, and Honggang Hu.
\newblock \emph{SQIsign2D2: Algorithm Specifications and Supporting Documentation}.
\newblock Version 1.0, NGCC submission, 26 April 2026.

\bibitem{sqisign2d2-eprint}
Zheng Xu, Kaizhan Lin, Chang-An Zhao, and Yi Ouyang.
\newblock SQIsign2D$^2$: New SQIsign2D Variant by Leveraging Power Smooth
Isogenies in Dimension One.
\newblock Cryptology ePrint Archive, Paper 2025/920, 2025.
\newblock \url{https://eprint.iacr.org/2025/920}.

\bibitem{ngcc-tracker}
Markku-Juhani O. Saarinen, editor.
\newblock SQIsign2D2 (sign-25), NGCC public security tracker.
\newblock Accessed 28 September 2026.
\newblock \url{https://ngcc.dev/reports/sign-25.html}.

\bibitem{fiat-shamir}
Amos Fiat and Adi Shamir.
\newblock How to Prove Yourself: Practical Solutions to Identification and Signature Problems.
\newblock In \emph{Advances in Cryptology -- CRYPTO '86}, LNCS 263, pages 186--194.
Springer, 1987.

\bibitem{gmr}
Shafi Goldwasser, Silvio Micali, and Ronald L. Rivest.
\newblock A Digital Signature Scheme Secure Against Adaptive Chosen-Message Attacks.
\newblock \emph{SIAM Journal on Computing}, 17(2):281--308, 1988.

\end{thebibliography}

\end{document}
